\section{Results}

\subsection{Ecological Validity and Data Quality}
Before interpreting physiological changes, we verified whether the two environments were indeed distinct in terms of signal quality. Figure 2 illustrates the distribution of the Exposure Fluctuation Index ($E_{fluc}$) across all valid windows.

As expected, the Home environment exhibited significantly higher baseline noise. The median $E_{fluc}$ in the Home context was $0.063$ (IQR = 0.041), nearly double that of the University context (Median = 0.032, IQR = 0.030). An independent samples t-test confirmed this difference to be highly significant ($t(196) = 7.15, p < 0.001$). This confirms that the "Home" dataset represents a true "in-the-wild" recording with greater environmental variability. Despite this increased noise, the Camerala tracking engine maintained a valid tracking rate of $>91\%$ in the Home condition, validating the robustness of the system.

\subsection{H1: The "Safe Haven" Interaction Effect}

Our central hypothesis (H1) postulated a **Context $\times$ Appraisal Interaction**: specifically, that the psychophysiological response to "Threat" framing would be significant in the University Lab but dampened or nonexistent in the Home environment.

\subsubsection{Behavioral Evidence (Reaction Time)}
We analyzed Reaction Time (RT) as a behavioral proxy for arousal and vigilance. In high-stakes environments, threat typically induces a "vigilance" response, accelerating RT (often at the cost of accuracy, though accuracy was clamped here).

We performed a 2-Way ANOVA (Context $\times$ Condition) on the RT data. The results revealed:
\begin{itemize}
    \item A significant Main Effect of Context ($F(1,396) = 9.30, p = 0.002$), indicating generally faster responses at Home.
    \item A significant Main Effect of Condition ($F(1,396) = 9.40, p = 0.002$).
    \item **Crucially, a significant Context $\times$ Condition Interaction** ($F(1,396) = 4.47, p = 0.035$).
\end{itemize}

To decompose this interaction, we analyzed the Threat vs. Challenge effect within each context:

\begin{itemize}
    \item \textbf{University Context}: "Threat" framing significantly accelerated Reaction Time ($M_{Threat}=767$ms) compared to "Challenge" ($M_{Chall}=1078$ms). The difference was statistically significant ($p = 0.004$), with a meaningful effect size ($d = -0.41$). Bayesian analysis yielded a Bayes Factor of **BF$_{10} = 66.7$**, constituting "Very Strong Evidence" for the effect.
    
    \item \textbf{Home Context}: The framing effect completely vanished. RT under Threat ($M=711$ms) was statistically indistinguishable from Challenge ($M=768$ms). The difference was non-significant ($p = 0.30$), and the effect size was negligible ($d = -0.15$). The Bayes Factor was **BF$_{10} = 1.7$**, indicating "Inconclusive/Anecdotal" evidence. 
\end{itemize}

This result (Figure 3) provides strong behavioral confirmation of the Safe Haven effect: the cognitive urgency induced by the "Threat" instruction was neutralized by the safety of the home.

\subsubsection{Physiological Evidence (Head Motion)}
We examined Head Motion ($M_t$) as a proxy for the "Freezing" response (reduced motion under threat).

The 2-Way ANOVA for Motion showed a trend toward interaction ($F(1,194)=1.71, p=0.19$), but the stratified analysis was revealing:
\begin{itemize}
    \item \textbf{University}: Threat induced a reduction in head motion (Freezing) compared to Challenge. This trend was **marginally significant** ($p = 0.08$) with a small effect size ($d = -0.33$). The Bayes Factor was **BF$_{10} = 4.6$** (Moderate Evidence).
    \item \textbf{Home}: Absolutely no difference was observed ($p = 0.86, d = -0.04$). The Bayes Factor was **BF$_{10} = 1.0$** (Null Evidence).
\end{itemize}

Table \ref{tab:effect_sizes} summarizes these contrasting effect sizes, clearly showing the degradation of the effect in the home environment.

\begin{table}[h]
\centering
\caption{Effect Size (Cohen's d) and Bayes Factor (BF10) Comparison}
\label{tab:effect_sizes}
\begin{tabular}{ll ccc}
\toprule
\textbf{Measure} & \textbf{Context} & \textbf{Cohen's d} & \textbf{BF$_{10}$} & \textbf{Interpretation} \\
\midrule
\multirow{2}{*}{Head Motion} & University & -0.33 & 4.6 & \textbf{Mod. Evidence} \\
                             & Home & -0.04 & 1.0 & Null \\
\midrule
\multirow{2}{*}{Reaction Time} & University & -0.41 & 66.7 & \textbf{Strong Evidence} \\
                              & Home & -0.15 & 1.7 & Inconclusive \\
\bottomrule
\end{tabular}
\end{table}

\subsection{H2: Physiological Invariance Breaking}
To test H2 (that physiological states are ambiguous without context), we aggregated all physiological features (Motion, Blink Rate, ROI Brightness) and performed Unsupervised K-Means Clustering ($k=3$).

If physiological arousal mapped 1:1 to affective state, we would expect one cluster to be predominantly "Threat" and another "Challenge." Instead, we found high entropy. Cluster 1 (the high-arousal cluster) contained a mixture of Threat (41\%) and Challenge (33\%) trials. This confirms that the physiological signal carries information about intensity, but lacks the valence/appraisal information necessary to distinguish Threat from Challenge. The "Context" variable is the missing key required to decode this ambiguity.

\subsection{Confirmatory Mixed-Effects Model}
Finally, to rule out the possibility that these results were driven by simple habituation (e.g., the participant just got used to the task in later sessions), we fitted a **Linear Mixed-Effects Model (LMM)** with `Session` as a random intercept using the `statsmodels` library in Python.

The model for RT confirmed that the **Context $\times$ Condition Interaction** remained significant ($z = -2.11, p = 0.034$) even after accounting for session-level variance. This demonstrates that the Safe Haven effect is a robust phenomenon, distinct from simple temporal habituation.
